\begin{abstract}
Nghiên cứu này đối chiếu hiệu quả của phản hồi từ công cụ độ phủ (JaCoCo) và công cụ đột biến (PITest) trong việc hỗ trợ sinh viên viết test case để bắt lỗi thực tế. Bằng thiết kế thực nghiệm chéo với 24 người tham gia, mỗi người được cấp 20 phút để viết thêm test JUnit cho hai chương trình thuộc IBM CodeNet, mỗi chương trình dùng một trong hai công cụ. Sự gia tăng trong tỉ lệ phát hiện các bài code lỗi (Wrong Answer) được giữ kín được dùng làm thang đo. Phân tích thống kê chưa thấy khác biệt giữa hai phương pháp (ước lượng Hodges--Lehmann $-0,012$, KTC 95\% $[-0,118; 0,074]$, Wilcoxon $p = 0,483$, $N = 22$). Thay vào đó, nghiên cứu phát hiện hai rào cản lớn: phần lớn test viết mới không hợp lệ, chủ yếu do sai giá trị kỳ vọng, và PITest với cấu hình mặc định dừng, không tạo báo cáo ở 58/96 lần chạy vì bộ test có test fail trên chương trình được kiểm thử. Kết quả bước đầu này gợi ý rằng, thay vì cân nhắc chọn loại báo cáo nào, các nhà giáo dục nên ưu tiên hướng dẫn sinh viên viết test hợp lệ và tinh chỉnh cấu hình PITest. Nghiên cứu thăm dò này cần được kiểm chứng lại với sinh viên thật và cỡ mẫu 48 như kế hoạch.
\end{abstract}
